% 389_Parabelflug_Darstellungen_De_ch.tex
% Johann Pascher, 6. Okt. 2026
% Parabelflug, Trägheit und die Darstellungen der Gravitation

\providecommand{\BM}{\textbf{[B]}}
\providecommand{\KM}{\textbf{[K]}}
\providecommand{\SM}{\textbf{[S]}}
\providecommand{\XM}{\textbf{[X]}}
\providecommand{\QM}{\textbf{[Q]}}

\chapter*{Parabelflug, Trägheit und die Darstellungen der Gravitation}
\addcontentsline{toc}{chapter}{Parabelflug, Trägheit und die Darstellungen der Gravitation}

\begin{abstract}
Im Parabelflug verschwindet die Schwerkraft für die Insassen. Dieses Dokument geht von dieser
Frage aus und liest sie mit der Zeit-Masse-Dualität $\tilde T\cdot m=1$: Der freie Fall ist die
Bahn stationärer Phase einer Mode mit ortsabhängiger Masse, und das Gewicht am Boden ist die
Abweichung der Ruhe-Weltlinie von dieser Bahn. Was sich nicht wegtransformieren lässt, sind die
Gezeiten. Gravitation erscheint dabei in drei Darstellungen, als Raumkrümmung, als Massen- bzw.\
Taktänderung und als fraktale Wegverlängerung. Sie sind mathematische Modelle, die aus der
Dualität heraus äquivalent sein müssen, keine Aussagen darüber, ob der Raum \glqq wirklich\grqq{}
flach oder gekrümmt ist. Für Licht sind Takt und Weg komplementäre Hälften: Jede allein ergibt
$0{,}875''$ am Sonnenrand, beide zusammen $1{,}75''$ und $\gamma=1$. Die rein konform flache
Metrik ergibt keine Ablenkung. Mit der logarithmischen Form der Massenänderung folgt auch
$\beta=1$ und damit die Periheldrehung der ART. Am Horizont begrenzt $T_0$ den Vergleich mit einem
fernen Bezug, nicht den freien Durchgang. Eine eigene Quantisierung der Gravitation wird nicht
gebraucht.
\end{abstract}

% ============================================================
\section*{Statusmarkierungen}
% ============================================================
\begin{center}
\begin{tabular}{@{}lp{11cm}@{}}
\textbf{[B]} & Algebraisch bewiesen. \\[2pt]
\textbf{[K]} & Numerisch bestätigt (Prüfskript, Vergleich mit Messwerten). \\[2pt]
\textbf{[S]} & Setzung, offen oder spekulativ. \\[2pt]
\textbf{[X]} & Widerlegt oder nicht tragfähig. \\[2pt]
\textbf{[Q]} & Konvention. \\
\end{tabular}
\end{center}
Messwerte: Pound--Rebka 1960, Cassini (Bertotti u.\,a.\ 2003), IAU- und CODATA-2022-Konstanten.
Prüfskript: \texttt{pruef\_389\_parabelflug.py}.

% ============================================================
\section{Die Frage}
% ============================================================

Im Parabelflug fliegt das Flugzeug für rund zwanzig Sekunden eine Wurfparabel. Die Insassen
schweben, obwohl die Erde unverändert anzieht. Am Boden dagegen spüren sie ihr Gewicht, obwohl
sie sich nicht bewegen. Die Frage ist, was dieses Gewicht ist, wenn es im freien Fall
verschwindet.

Die Antwort, die dieses Dokument ausarbeitet: Die Kraft, die wir als Gewicht spüren, ist nur
vorhanden, weil wir von der Bahn abweichen, die wir ohne Unterlage nehmen würden. Der Boden zwingt
uns auf eine Bahn, die der Fallbahn genau entgegengesetzt beschleunigt ist. Im Parabelflug fällt
diese Abweichung weg, und mit ihr das Gewicht. Die folgenden Abschnitte rechnen das in der
Darstellung der FFGFT nach und ordnen es in die verschiedenen Lesarten der Gravitation ein.

% ============================================================
\section{Freier Fall in der Massendarstellung}
% ============================================================

In der Darstellung \glqq Raum flach, Masse veränderlich\grqq{} (Dok.~078) trägt eine Mode im
Potential $\Phi$ die ortsabhängige Masse $m(x)$. Ihre Wirkung ist die einer Uhr, die mit der
Rate $m(x)c^2$ tickt:
\begin{equation}
L=-m(x)\,c^2\sqrt{1-v^2/c^2},\qquad m(x)=m_0\,e^{\Phi(x)/c^2}.
\end{equation}
Dabei ist $m(x)$ die vom fernen Bezug aus verglichene Masse. Eine Uhr bei $\Phi$ tickt in dessen
Zeit mit
\begin{equation}
\omega(\Phi)=\omega_0\sqrt{g_{00}}\approx\omega_0\Bigl(1+\frac{\Phi}{c^2}\Bigr),\qquad
m(x)=m_0\sqrt{g_{00}},\qquad g_{00}=\Bigl(\frac{m}{m_0}\Bigr)^2=e^{2\Phi/c^2}\quad\BM,
\end{equation}
denn mit $\tilde T\cdot m=1$ ist die Masse die Frequenz der Mode. $m(x)$ ist die Killing-Energie
der Mode; mit einer mitgeführten Uhr gemessen bleibt die Masse $m_0$.
Die Euler-Lagrange-Gleichung gibt für kleine Geschwindigkeit
\begin{equation}
\ddot{\vec x}=-c^2\,\nabla\ln m=c^2\,\nabla\ln\lambda_4=-\nabla\Phi\quad\BM,
\end{equation}
mit der Modenlänge $\lambda_4=2\pi/m$ (Dok.~314). Ohne Metrik und ohne Christoffelsymbole folgt der
Fall allein aus dem Gradienten der Masse: Körper fallen dorthin, wo ihre Masse kleiner und ihre
Modenlänge größer ist. Weil $\nabla\ln m$ für jede Mode derselbe ist, fallen alle Körper gleich;
das Äquivalenzprinzip ist in dieser Form eingebaut.

An der Erdoberfläche ist $g/c^2=1{,}09\times10^{-16}$ pro Meter \KM. Das ist dieselbe Zahl, um die
eine Uhr pro Meter Höhe schneller geht. Fallbeschleunigung und Uhrengang sind in dieser
Darstellung eine einzige Größe.

\paragraph{Vorzeichen.} Die Richtung der Massenänderung ist durch die Rotverschiebung festgelegt.
Licht, das im Schwerefeld aufsteigt, kommt rotverschoben an; der Pound-Rebka-Versuch hat das über
$22{,}5$\,m mit $(2{,}57\pm0{,}26)\times10^{-15}$ gemessen, erwartet sind $2{,}46\times10^{-15}$
\KM. In der Massendarstellung heißt das: Die Mode unten hat, vom oberen Bezug aus verglichen, die
kleinere Masse, $m$ wächst mit $\Phi$. Mit diesem Vorzeichen fällt der Körper nach unten. Das
Vorzeichen ist damit keine offene Wahl, sondern durch die Messung bestimmt.

% ============================================================
\section{Ruhe als Gegenbahn}
% ============================================================

Die Phase einer Mode entlang eines Weges ist $\varphi=-\int m(x)\,c^2\,d\tau$. Der freie Fall ist
die Bahn, auf der diese Phase stationär ist, $\delta\varphi=0$. Die Ruhe-Weltlinie
$\vec x=\text{const}$ ist es nicht: Ihre Variation verschwindet nicht, sondern lässt den Rest
\begin{equation}
\frac{\partial L}{\partial \vec x}\Big|_{v=0}=-m\,\nabla\Phi=-m\,\vec g\quad\BM.
\end{equation}
Damit ein Körper trotzdem ruht, muss der Boden genau diesen Rest mit $+m\vec g$ nach oben
ausgleichen. Das ist das Gewicht. Es ist keine zusätzliche Kraft, die zur Schwerkraft hinzukommt,
sondern das Maß dafür, wie weit die Ruhe-Weltlinie von der Bahn stationärer Phase abweicht. In
der Sprache der Mode ist es eine Phasenverstimmung: Die ruhende Mode läuft gegen den Takt, den sie
im freien Fall hätte.

Im Parabelflug folgt das Flugzeug der Fallbahn. Die Abweichung verschwindet, und mit ihr das
Gewicht. Was die Insassen am Boden als Schwere und im Flugzeug als Schwerelosigkeit erleben, ist
dieselbe Bahn, einmal mit und einmal ohne Gegenbahn. Als Formel:
\begin{equation}
\vec a_{\text{Fall}}=c^2\,\nabla\ln\lambda_4,\qquad
\vec a_{\text{Träg}}=-c^2\,\nabla\ln\lambda_4=-\vec a_{\text{Fall}}\quad\BM.
\end{equation}
Am Boden heben sich beide auf, im Parabelflug ist die Gegenbahn null. In diesem Sinn ist die
Trägheitskraft nur vorhanden, weil wir eine genau entgegengesetzte Bahn haben.

% ============================================================
\section{Was sich nicht wegtransformieren lässt}
% ============================================================

Der Gradient $\nabla\ln m$ lässt sich lokal wegtransformieren: Im mitfallenden System ist er null,
wie in der ART die Christoffelsymbole im lokalen Inertialsystem. Die zweite Ableitung lässt sich
nicht wegtransformieren. Die Gezeiten sind die Hesse-Matrix
\begin{equation}
\partial_i\partial_j\Phi=c^2\,\partial_i\partial_j\ln m,\qquad
\text{Eigenwerte }(-2,\,1,\,1)\,\frac{GM}{r^3}\ \text{auf der Achse},\qquad
\nabla^2\ln m=\frac{4\pi G\rho}{c^2}\quad\BM.
\end{equation}
Im Vakuum ist die Spur null, im Inneren einer Masse gibt sie die Poisson-Gleichung, jetzt als
Gleichung für das Massenfeld. Außen ist $m=m_0\,e^{-GM/(rc^2)}\approx m_0\bigl(1-GM/(rc^2)\bigr)$;
der Gradient $\partial_r\ln m=GM/(r^2c^2)$ lässt sich wegtransformieren, die zweite Ableitung
$c^2\,\partial_r^2\ln m=-2GM/r^3$ nicht \BM. An der Erdoberfläche ist $2GM/r^3=3{,}08\times10^{-6}\,\text{s}^{-2}$
\KM. Auch im Parabelflug bleibt dieser Rest: Zwei Körper an verschiedenen Stellen des Flugzeugs
fallen nicht exakt gleich. Was die Massendarstellung \glqq Krümmung\grqq{} nennt, ist die Krümmung
des Massenfeldes, nicht die einer Raumgeometrie.

% ============================================================
\section{Drei Darstellungen}
% ============================================================

Eine ortsabhängige Masse ist gleichwertig mit einer ortsabhängigen Zeit, denn
$\tilde T\cdot m=1$. Wie man das nennt, ist eine Wahl der Darstellung:
\begin{itemize}[leftmargin=*,itemsep=1pt]
\item \textbf{Raumkrümmung:} Masse konstant, Raum und Zeit gekrümmt, wie in der ART.
\item \textbf{Massen- bzw.\ Taktänderung:} Raum flach, Masse und Takt ortsabhängig (Dok.~078).
\item \textbf{Fraktale Wegverlängerung:} Raum flach, Ausbreitungswege geometrisch gestreckt
(Dok.~026, 041, 182).
\end{itemize}
Es sind mathematische Modelle, die aus der Dualität heraus äquivalent sein müssen. Ontologisch
lässt sich daraus weder sagen, der Raum sei flach, noch, er sei gekrümmt; Mischformen sind
denkbar. Entscheidend ist allein, dass die Observablen übereinstimmen.

\paragraph{Licht als Prüfstein.} Für Uhren, Rotverschiebung und langsame Körper genügt der Takt.
Licht prüft mehr, denn es spürt zusätzlich die räumliche Geometrie. In der Fermat-Form läuft Licht
durch einen Brechungsindex $n=1+(1+\gamma)\,GM/(rc^2)$; die Ablenkung am Sonnenrand und die
Shapiro-Verzögerung skalieren mit $(1+\gamma)/2$:
\begin{center}
\begin{tabular}{lcc}
\toprule
Darstellung & $\gamma$ & Ablenkung am Sonnenrand \\
\midrule
rein konform flache Metrik (Dok.~143) & $-1$ & $0$ \\
nur Takt (Einstein 1911) & $0$ & $0{,}875''$ \\
Takt und fraktale Wegverlängerung (Dok.~308) & $1$ & $1{,}751''$ \\
\bottomrule
\end{tabular}
\end{center}
\KM. Cassini misst $\gamma-1=(2{,}1\pm2{,}3)\times10^{-5}$; nur die letzte Zeile ist mit der
Messung verträglich. Dok.~308 rechnet dasselbe als Strahlverfolgung durch: Takt und
Wegverlängerung ergeben je $0{,}875''$, zusammen $1{,}750''$. Sie sind komplementäre Hälften,
keine Alternativen.

Daraus folgt eine Präzisierung der drei Darstellungen. \glqq Raum flach, Masse veränderlich\grqq{}
allein trägt nur die zeitliche Hälfte. Vollständig wird die flache Darstellung erst mit der
fraktalen Wegverlängerung, die die räumliche Hälfte trägt. Die Raumkrümmung der ART ist dann
gleichwertig mit Takt und Weg zusammen. Die zeitliche Hälfte ist durch Uhren gemessen und in jeder
Darstellung dieselbe; die räumliche Hälfte lässt sich als Krümmung, als Wegverlängerung oder als
Mischung aus beiden schreiben. Die Freiheit der Darstellung liegt also in der räumlichen Hälfte,
und die Bedingung für jede Mischform ist $\gamma=1$.

Die konform flache Metrik
$g_{\mu\nu}=\eta_{\mu\nu}(1-2\xi E_{\text{field}}/E_P)$ aus Dok.~143 ist keine dieser Darstellungen:
Sie streckt Raum und Zeit mit demselben Faktor, Nullgeodäten bleiben dabei unverändert, und Licht
läuft wie im flachen Raum. Das ist der Nordström-Fall (Dok.~143). Dasselbe gilt für die konforme Kopplung $g_{\mu\nu}\to\Omega^2g_{\mu\nu}$ der
Materiefelder an das Zeitfeld (A142): Das Maxwell-Feld ist in vier Dimensionen konform invariant,
die Kopplung lässt Licht unberührt \BM. Für massive Materie liefert sie den Takt, also die
zeitliche Hälfte; die räumliche Hälfte muss aus der Wegverlängerung kommen (A142).

% ============================================================
\section{Zweite Ordnung: die Periheldrehung}
% ============================================================

Die Lichtablenkung prüft die erste Ordnung im Potential. Die Periheldrehung prüft zusätzlich die
zweite Ordnung der Zeitkomponente, den Parameter $\beta$ in $g_{00}=1-2U+2\beta U^2$ mit
$U=GM/(rc^2)$. Der Faktor gegenüber der ART ist $(2+2\gamma-\beta)/3$. In der Massendarstellung ist
$g_{00}=(m/m_0)^2$, und $\beta$ hängt an der Form von $m(\Phi)$:
\begin{center}
\begin{tabular}{lccc}
\toprule
Form & $\beta$ & mit $\gamma$ & Merkur ($''$/Jh.) \\
\midrule
$m=m_0\,e^{\Phi/c^2}$ & $1$ & $1$ & $42{,}98$ (ART) \\
$m=m_0\,(1+\Phi/c^2)$ & $\tfrac12$ & $1$ & $50{,}1$ \\
Nordström & $\tfrac12$ & $-1$ & $-7{,}2$ \\
nur Takt & $1$ & $0$ & $14{,}3$ \\
\bottomrule
\end{tabular}
\end{center}
\KM. Die logarithmische Form $\ln(m/m_0)=\Phi/c^2$ ist dieselbe, die in Abschnitt~2 den Fall
$\ddot{\vec x}=-\nabla\Phi$ exakt und nicht nur genähert liefert. Mit ihr und mit der räumlichen
Hälfte aus Dok.~308 gibt die flache Darstellung die Periheldrehung der ART. Die logarithmische
Form selbst ist gesetzt, nicht aus $\tilde T\cdot m=1$ hergeleitet \SM; die lineare Form ist
durch die Merkur-Messung ausgeschlossen.

% ============================================================
\section{Die Relation \texorpdfstring{$\lambda_4\cdot m=2\pi$}{lambda4 m = 2 pi} gilt lokal}
% ============================================================

Die Ortsrelation $\lambda_4\cdot m=2\pi$ (Dok.~314) verbindet Eigengrößen: die Modenlänge, wie
ein mitgeführter Maßstab sie misst, und die Masse, wie eine mitgeführte Uhr sie misst. In
Koordinaten des fernen Bezugs gilt sie nicht unverändert, sobald die räumliche Hälfte als
Krümmung geschrieben wird, denn dann tragen Koordinatenlängen einen Faktor der räumlichen Metrik.
In der isotropen Form der ART ist
\begin{equation}
g_{ij}=\Bigl(1-\frac{2\Phi}{c^2}\Bigr)\delta_{ij},\qquad
\ell_{\text{koord}}=\ell\,\Bigl(1+\frac{\Phi}{c^2}\Bigr),\qquad
\lambda_{4,\text{koord}}=\frac{2\pi}{m(x)}=\lambda_4\Bigl(1-\frac{\Phi}{c^2}\Bigr)\quad\BM.
\end{equation}
Räumliche Längen erscheinen vom fernen Bezug aus tief im Potential verkürzt, die zeitliche
Modenlänge verlängert: Die beiden Faktoren laufen gegeneinander. In der rein konform flachen
Metrik liefen beide mit $1/m$ gleich, und Licht bliebe unabgelenkt. Der räumliche Faktor ist genau
die Hälfte, die Licht zusätzlich spürt. In der Darstellung mit fraktaler Wegverlängerung steckt
er im gestreckten Weg, als Brechungsanteil $n-1=-\Phi/c^2$ zum Takt hinzu. Der
Gradient $\nabla\ln\lambda_4$ in Abschnitt~2 ist entsprechend eine lokale Größe: Im mitfallenden
System verschwindet er, im ruhenden ist er $g/c^2$.

% ============================================================
\section{Horizont: Vergleich und Durchgang}
% ============================================================

Dok.~306 beschreibt den Horizont als Extremum: Dort erreicht $T$ den Wert $T_0=\xi\,t_P$ und die
Energie je Anregung den Wert $E_P/\xi$. Es fügt hinzu, dass die \glqq langsamste Zeit\grqq{} am
Horizont eine Aussage über den Vergleich ist, nicht über das Erleben. Die Rechnung bestätigt das.

Ein statischer Beobachter bei $r$ misst für eine Mode der Energie $E$ (bezogen auf den fernen
Bezug) die lokale Energie $E/\sqrt{1-r_s/r}$; sie divergiert am Horizont \BM. Die Schranke
$E\le E_P/\xi$ schneidet diese Annäherung ab, beim Eigenabstand
$\rho_{\min}=2\,r_s\,\xi\,E/E_P$ \BM. Ein frei fallender Beobachter dagegen misst für einlaufendes
Licht $E/(1+\sqrt{r_s/r})$, am Horizont also $E/2$, endlich \BM. Für ihn ist am Horizont nichts
zu begrenzen.

Die Schranke $T_0$ betrifft damit die statische Lesart, den Vergleich mit einem fernen Bezug.
Der frei fallende Durchgang ist davon nicht berührt; das ist dieselbe Unterscheidung wie zwischen
Ruhe und Fall im Parabelflug. Ob $T_0$ darüber hinaus den Radius im Inneren nach unten begrenzt,
verlangt eine Dynamik des Inneren, die der Korpus nicht liefert. Die physikalische Gültigkeit im
Regime Schwarzer Löcher bleibt offen \SM\ (Dok.~384).

% ============================================================
\section{Zwei Massen, zwei Lesarten von \texorpdfstring{$G$}{G}}
% ============================================================

Zwei Massen am oberen Ende werden leicht verwechselt:
\begin{itemize}[leftmargin=*,itemsep=1pt]
\item $M_{\text{coll}}=m_P/\sqrt2$: Bei ihr ist der Schwarzschild-Radius gleich der reduzierten
Compton-Länge \BM\ (Dok.~325, 329).
\item $m_P/\xi=7500\,m_P\approx1{,}63\times10^{-4}$\,kg: die duale Obergrenze je Anregung,
zu $L_0=\xi\,\ell_P$ gehörig \KM\ (Dok.~290, 306).
\end{itemize}
Beide sind richtig, aber verschieden; ihr Verhältnis ist $7500\sqrt2\approx10\,607$. Nach
Dok.~306 sind beide Schranken je Anregung, nicht je Aggregat: Ein Stern ist keine einzelne
Windung, und $\hbar/(Mc^2)$ für eine Sternmasse ist kein Zustand der Relation.

Für $G$ gilt Entsprechendes. Mit dem gemessenen, festen $G$ gilt $m\cdot G=\xi^2/4$ für die
charakteristische Masse $m_{\text{char}}$ des Ankers (R58, A142). Für beliebige $m$ gilt die
Relation nur, wenn $G$ mitläuft, wie Dok.~350 es für ein verdampfendes Loch liest. Die Aussage
\glqq keine Massensingularität mit endlichem $G$\grqq{} ist in der ersten Lesart eine Folge der
Festlegung, in der zweiten eine Aussage über die Dynamik.

% ============================================================
\section{Keine eigene Quantisierung der Gravitation}
% ============================================================

In der FFGFT ist Gravitation kein zusätzliches Feld und braucht kein Graviton. Sie ist in der
Zeitfeld-Geometrie bereits enthalten und zeigt sich als Wirkung in der Rotverschiebung, gelesen
als Raumkrümmung, als Massenänderung oder als fraktale Wegverlängerung. Eine eigene Quantisierung
der Gravitation wird deshalb nicht gefordert (R58, A142). Offen bleibt die Überführung der
klassischen Gleichungen in Schleifenkorrekturen für Situationen jenseits des
Schwarzschild-Grenzfalls (A142).

% ============================================================
\section{Was offen bleibt}
% ============================================================

\begin{itemize}[leftmargin=*,itemsep=1pt]
\item die logarithmische Form $\ln(m/m_0)=\Phi/c^2$ als Folge der Dualität, nicht nur als
Setzung, die $\beta=1$ liefert \SM;
\item die Herkunft der fraktalen Wegverlängerung als räumliche Hälfte aus der Struktur des
Trägers; Dok.~308 führt die K3-Flächeninvarianz als Begründung, numerisch bestätigt, aber nicht
bewiesen \KM;
\item die Dynamik im Inneren eines Horizonts und die Gültigkeit von $m\cdot G=\xi^2/4$ jenseits
des Ankers \SM.
\end{itemize}

% ============================================================
\section{Fazit}
% ============================================================

Im Parabelflug verschwindet nicht die Gravitation, sondern die Gegenbahn. Der freie Fall ist die
Bahn stationärer Phase einer Mode mit ortsabhängiger Masse, das Gewicht am Boden ist die
Abweichung der Ruhe-Weltlinie von dieser Bahn, und nur die Gezeiten bleiben in jedem System
bestehen. Raumkrümmung, Massenänderung und fraktale Wegverlängerung sind gleichwertige Modelle
desselben Sachverhalts, keine Aussagen über das Wesen des Raumes. Für Licht sind Takt und Weg
komplementäre Hälften; erst beide zusammen geben $1{,}75''$ und $\gamma=1$. Mit der
logarithmischen Form der Massenänderung folgt auch die Periheldrehung der ART. Am Horizont trennt
dieselbe Unterscheidung zwischen Ruhe und Fall den Vergleich mit einem fernen Bezug vom freien
Durchgang.

\vspace{1em}
\noindent\textit{Prüfskript:} \texttt{2/python/Dok389\_Skripte/pruef\_389\_parabelflug.py}
--- 34/34.
